% TOP_PROBLEM: {"id": 4600003, "title": "Range of the dimension representation", "category_id": 11, "category": {"id": 11, "name": "dynamical_systems", "display_name": "Dynamical Systems", "description": "Problems about long-term behavior of deterministic systems, Hamiltonian dynamics, and periodic orbits.", "slug": "dynamical-systems", "order_index": 11, "created_at": "2026-07-31T15:26:25.671Z"}, "status": "open", "problem_number": "AMR-045-0003", "set_id": 13}
% FIRST_POSTED: 2026-10-02
% ENTRY_KIND: statement_only
% UnsolvedMath ID 4600003; source catalogue code AMR-045-0003
% !TeX program = lualatex
% Definition and sources only; this file is not an LLM solution attempt.
\documentclass[11pt,a4paper]{article}
\usepackage[margin=25mm]{geometry}
\usepackage{fontspec}
\setmainfont{lmroman10-regular.otf}[BoldFont=lmroman10-bold.otf,
 ItalicFont=lmroman10-italic.otf,BoldItalicFont=lmroman10-bolditalic.otf]
\usepackage{amsmath,amssymb,mathtools}
\usepackage{unicode-math}
\setmathfont{latinmodern-math.otf}
\AtBeginDocument{\let\setminus\smallsetminus}
\usepackage{xurl,hyperref}
\hypersetup{colorlinks=true,urlcolor=blue}
\setcounter{secnumdepth}{0}
\setlength{\parindent}{0pt}
\setlength{\parskip}{5pt}
\setlength{\emergencystretch}{3em}
\begin{document}
\section*{Range of the dimension representation}\label{4600003}
{\small UnsolvedMath ID 4600003 · AMR-045-0003 · Dynamical Systems · AMR Open Problem Lists}
\subsection{Definitions and mathematical statement}
Let $A$ be a primitive square matrix with nonnegative integer entries.
The edge shift $X_A$ consists of bi-infinite paths in its directed
multigraph, with shift $\sigma$ deleting the initial coordinate. Primitivity
means that some power of $A$ is strictly positive and ensures mixing.
Use row vectors and define the dimension group
\[
 G_A=\varinjlim(\mathbb Z^N,\ v\mapsto vA),
 \qquad [v,j]=[vA,j+1].
\]
Its positive cone consists of classes represented by a nonnegative
integer vector at some stage. The dimension automorphism is
$\delta_A[v,j]=[vA,j]$; its inverse is $[v,j]\mapsto[v,j+1]$.

Every topological automorphism $\phi$ of $(X_A,\sigma)$ induces an
order-preserving automorphism $\rho_A(\phi)$ of $G_A$ commuting with
$\delta_A$. To specify this representation concretely, a strong shift
equivalence step $A=UV$, $B=VU$ induces the map
$[v,j]\mapsto[vU,j]$. A chain of such nonnegative integral steps, followed
by a shift power, represents any conjugacy and defines its induced map;
the result is independent of that representation.

The problem is to determine exactly the subgroup
\[
 \rho_A(\operatorname{Aut}(X_A,\sigma))
       \subseteq\operatorname{Aut}(G_A,G_A^+,\delta_A).
\]
Thus one seeks necessary and sufficient conditions for a dimension-module
automorphism to lift to a sliding-block automorphism of the shift.
Merely describing all algebraic automorphisms of the ordered group does
not settle liftability. The matrix $A$ is arbitrary, and both the positive
cone and the commuting condition are part of the target structure.
\subsection{Short English statement}
Determine which automorphisms of a mixing shift’s ordered dimension module come from actual shift automorphisms.
\subsection{Sources}
\begin{itemize}
\item[S1] ULAM.AI and the UnsolvedMath Contributors, supplied problems.json, record 4600003 (AMR-045-0003), AMR Open Problem Lists. Catalogue identity and source transcription; CC BY 4.0.
\url{https://huggingface.co/datasets/ulamai/UnsolvedMath}.
\item[S2] Boyle - Open Problems in Symbolic Dynamics (2008), Problem/Question/Conjecture 4.1. Original source indexed by AMR-045-0003.
\url{https://www.math.umd.edu/~mboyle/papers/openfinalsub3nov2007.pdf}.
\end{itemize}
\end{document}
